\documentclass[10pt,a4paper]{amsart}
\usepackage[margin=19mm]{geometry}
\usepackage{amsmath,amssymb,mathtools}
\usepackage[scr=rsfs,cal=euler]{mathalfa}
\usepackage{xfrac}
\usepackage[unicode,hidelinks,bookmarks=false]{hyperref}
\newcommand{\ie}{i.e.\ }
\newcommand{\cf}{cf.\ }
\newcommand{\resp}{respectively\ }
\newcommand{\Subsection}{\subsection}
\providecommand{\nobreakdash}{\nobreak\hbox{-}}
\setlength{\emergencystretch}{2em}
\multlinegap0pt
\usepackage{amssymb}
\newtheorem{lemma}{Lemma}
\def\bN{\mathbb{N}}
\def\fc{{\mathfrak{c}}}
\def\fd{{\mathfrak{d}}}
\def\fe{{\mathfrak{e}}}
\title{\textbf{A complete ultrametric on von Neumann's incomplete tensor products}}
\author{Andrew Lesniewski}
\date{Source: arXiv:2607.09627v3 (CC BY 4.0)}
\begin{document}
\maketitle

\noindent\emph{Source and licence.} \textbf{A complete ultrametric on von Neumann's incomplete tensor products}. Version arXiv:2607.09627v3 (CC BY 4.0), distributed by arXiv under the Creative Commons Attribution 4.0 International licence, http://creativecommons.org/licenses/by/4.0/, as recorded in the arXiv licence metadata for this exact version and shown on the abstract page https://arxiv.org/abs/2607.09627v3. Full licence notice: \texttt{licenses/arxiv-2607.09627-CC-BY-4.0.txt}.\par\smallskip

\noindent\textit{Editorial scope.} Counted unit: the article's lemma complLem (source0.tex lines 461--463). The article's complete original proof of it is delivered with it (source0.tex lines 464--530, one \texttt{proof} environment of 67 lines). No mathematics is written, repaired or re-derived: the statement and the proof are the article's own lines, in the article's own notation, and no retained line is substituted or re-flowed.  Retained uncounted as the source-local context the proof needs: lines 218--220; lines 292--294; lines 299--301; lines 350--352; lines 403--410.  Nothing was omitted from any retained span, so the package carries no omission record.  No retained line contains a citation, so the wrapper carries no bibliography and no bibliography entry.  No retained line contains a figure environment, an image-inclusion command or drawing code, so no image is delivered and the package declares no asset dependency.  Editorial formatting only, in addition: an AMS \texttt{amsart} preamble that restates the article's own declarations and macros used by the retained lines, namely the environments \texttt{lemma} on separate counters, together with the standard packages those lines need.  The retained environments print in this document as Lemma 1 in order of appearance, whereas the article prints the same items with the numbering of its own full document.  The complete native source of the article as distributed in the arXiv e-print for this exact version is bundled unchanged as \texttt{sources/204883/source0.tex}.
\begin{lemma}\label{unitRep}
Each equivalence class $\fc$ has a representative $\varphi\in\fc$ such that $\|\varphi_\iota\|=1$, for all $\iota$.
\end{lemma}

\begin{equation}\label{dDef}
d(\varphi,\psi)=\inf\Big\{p\geq 0:\:\sum _{j=1}^\infty\,\frac{1}{j^p}\,|\langle \varphi_j,\psi_j\rangle-1|<\infty\Big\}.
\end{equation}

\begin{equation}\label{upSet}
p>d(\varphi,\psi)\quad\Longrightarrow\quad\sum _{j=1}^\infty\,\frac{1}{j^p}\,|\langle \varphi_j,\psi_j\rangle-1|<\infty,
\end{equation}

\begin{lemma}\label{classInv}
If $\varphi\sim\varphi'$, then $d(\varphi,\psi)=d(\varphi',\psi)$. Consequently, \eqref{dDef} defines a function $d(\fc,\fd)$ on $\Gamma\times\Gamma$.
\end{lemma}

\begin{lemma}\label{ultraLem}
The function $d$ defines a pseudo-ultrametric on $\Gamma$, i.e.\ it has the following properties:
\begin{itemize}
\item[(i)]{$d(\fc,\fd)\geq 0$, and $d(\fc,\fc)=0$;}
\item[(ii)]{for all $\fc,\fd\in\Gamma$, $d(\fc,\fd)=d(\fd,\fc)$;}
\item[(iii)]{for all $\fc,\fd,\fe\in\Gamma$, $d(\fc,\fe)\leq \max\big(d(\fc,\fd),d(\fd,\fe)\big)$.}
\end{itemize}
\end{lemma}

\begin{lemma}\label{complLem}
Any Cauchy sequence in $(\Gamma,d)$, $\fc_n=[\varphi_n]$, $n=1,2,\ldots$, has a limit.
\end{lemma}

\begin{proof}
By Lemma \ref{unitRep} and Lemma \ref{classInv}, we may and do choose the representatives $\varphi_n\in\fc_n$ so that $\|\varphi_{n,j}\|=1$, for all $n$ and $j$. Note that then
\begin{equation}\label{unitBound}
|\langle\varphi_{n,j},\varphi_{m,j}\rangle-1|\leq 2,\qquad\text{for all }n,m,j.
\end{equation}

\smallskip\noindent\textit{Step 1 (a rapidly Cauchy subsequence).}
Since $(\fc_n)$ is Cauchy, we may choose indices $n_1<n_2<\ldots$ such that
\begin{equation}\label{rapidCauchy}
d(\fc_m,\fc_{m'})<2^{-k},
\end{equation}
for all $m,m'\geq n_k$. In particular, $d(\varphi_{n_k},\varphi_{n_l})<2^{-k}$ for all $l>k$, and therefore, by \eqref{upSet},
\begin{equation}\label{pairSeries}
\sum_{j=1}^\infty\,\frac{1}{j^p}\,|\langle \varphi_{n_k,j},\varphi_{n_l,j}\rangle-1|<\infty,
\end{equation}
for every $p>2^{-k}$ and every $l>k$. We emphasize that \eqref{pairSeries} asserts convergence for each fixed pair $(k,l)$; no uniformity over the pairs is claimed, and none will be needed.

\smallskip\noindent\textit{Step 2 (choice of blocks).}
For $l\in\bN$, introduce the finite set of test exponents
\begin{equation*}
P_l=\Big\{2^{-k}\Big(1+\frac{1}{m}\Big):\:1\leq k\leq l,\ 1\leq m\leq l\Big\}.
\end{equation*}
We define integers $1=J_1<J_2<J_3<\ldots$ recursively, as follows. Suppose $J_1,\ldots,J_{l-1}$ have been chosen. For every $k<l$ and every $p\in P_l$ with $p>2^{-k}$, the series in \eqref{pairSeries} converges, and hence its tails tend to zero. Since there are only finitely many such pairs $(k,p)$, we may choose $J_l>J_{l-1}$ so large that
\begin{equation}\label{tailBound}
\sum_{j\geq J_l}\,\frac{1}{j^p}\,|\langle \varphi_{n_k,j},\varphi_{n_l,j}\rangle-1|\leq 2^{-l},
\end{equation}
for all $1\leq k<l$ and all $p\in P_l$, with $p>2^{-k}.$ Define the blocks $B_l=\{j\in\bN:\:J_l\leq j<J_{l+1}\}$, $l=1,2,\ldots$; they partition $\bN$. Now set
\begin{equation*}
\varphi^\ast_j=\varphi_{n_l,j},\qquad\text{for }j\in B_l .
\end{equation*}
Since $\|\varphi^\ast_j\|=1$ for all $j$, the sequence $\varphi^\ast$ is a $C_0$-sequence; let $\fc^\ast=[\varphi^\ast]\in\Gamma$.

\smallskip\noindent\textit{Step 3 (convergence along the subsequence).}
We claim that
\begin{equation}\label{subConv}
d(\fc^\ast,\fc_{n_k})\leq 2^{-k},\qquad\text{for every }k.
\end{equation}
Fix $k$ and let $p>2^{-k}$ be arbitrary. Choose an integer $m\geq1$ with
\begin{equation*}
p'=2^{-k}\Big(1+\frac1m\Big)\leq p,
\end{equation*}
and set $l_0=\max(k+1,m)$. Splitting the series over the blocks $B_l$, we obtain
\begin{equation*}
\sum_{j=1}^\infty\,\frac{1}{j^{p'}}\,|\langle \varphi^\ast_{j},\varphi_{n_k,j}\rangle-1|
=\sum_{l=1}^\infty\;\sum_{j\in B_l}\,\frac{1}{j^{p'}}\,|\langle \varphi_{n_l,j},\varphi_{n_k,j}\rangle-1|.
\end{equation*}
The blocks with $l<l_0$ are finite sets, and their contribution is finite by \eqref{unitBound}. For $l\geq l_0$ we have $k<l$, $p'\in P_l$, and $p'>2^{-k}$, and so, by \eqref{tailBound},
\begin{equation*}
\begin{split}
\sum_{j\in B_l}\,\frac{1}{j^{p'}}\,|\langle \varphi_{n_l,j},\varphi_{n_k,j}\rangle-1|
&\leq\sum_{j\geq J_l}\,\frac{1}{j^{p'}}\,|\langle \varphi_{n_k,j},\varphi_{n_l,j}\rangle-1|\\
&\leq 2^{-l}.
\end{split}
\end{equation*}
Summing over $l\geq l_0$ yields a finite contribution as well. Consequently, the full series converges at the exponent $p'$, and thus $d(\fc^\ast,\fc_{n_k})\leq p'\leq p$. Since $p>2^{-k}$ was arbitrary, \eqref{subConv} follows.

\smallskip\noindent\textit{Step 4 (convergence of the full sequence).}
Let $\epsilon>0$, and choose $k$ with $2^{-k}<\epsilon$. For $n\geq n_k$, the strong triangle inequality (Lemma \ref{ultraLem} (iii)), together with \eqref{subConv} and \eqref{rapidCauchy}, gives
\begin{equation*}
\begin{split}
d(\fc^\ast,\fc_n)&\leq\max\big(d(\fc^\ast,\fc_{n_k}),\,d(\fc_{n_k},\fc_n)\big)\\
&\leq 2^{-k}\\
&<\epsilon.
\end{split}
\end{equation*}
Hence $\lim_{n\to\infty}d(\fc^\ast,\fc_n)=0$, and $\fc^\ast$ is a limit of the sequence $(\fc_n)$.
\end{proof}

\end{document}
